% =============================================================================
% Introduction, version 4 (27 November 2026)
% Light polish of the author-supplied five-paragraph draft, with citations added.
% Drop-in replacement for the \section*{Introduction} block of
% "Manuscript_full .tex". Citation keys resolve against references_2.bib plus
% Introduction_V3_new_refs.bib (panchanathan2004, rankin2009).
% \cite is used without a preceding tilde, as in Manuscript_full .tex, because
% the class loads cite with the superscript option (a tilde would leave a gap).
% British spelling, as in the supplied draft.
% Citation map and checks: manuscript_planning/Introduction_V4_notes.md
% =============================================================================
\section*{Introduction}

Reputation changes the evolutionary consequences of social behaviour. By conditioning cooperation on the reputations of others, individuals can reward those who have cooperated in the past, even if they have never interacted directly. This mechanism, known as indirect reciprocity, allows individuals to recover the costs of cooperation through benefits received from third parties rather than from the original recipients\cite{alexander1987,nowak1998,nowak2005}. Since the seminal work on image scoring\cite{nowak1998}, theoretical and experimental studies have established the role of indirect reciprocity in the evolution of cooperation, from pairwise interactions\cite{wedekind2000,leimar2001,schmid2021nhb} to collective action\cite{milinski2002,panchanathan2004,yoeli2013}. In this process, reputations serve two closely related functions: they provide social information that guides cooperative decisions and generate evolutionary incentives for cooperation\cite{milinski2016,nowak2005}. Both functions depend on the ability of individuals to form reliable judgements about others. How reputations are formed and maintained is therefore central to understanding how indirect reciprocity sustains cooperation.

Much of the theoretical research on indirect reciprocity has focused on identifying social norms under which cooperation is evolutionarily stable, including the ``leading eight''\cite{brandt2004,ohtsuki2004,pacheco2006,ohtsuki2006,santos2018,kessinger2023,murase2024computational}. A key insight from this research is that maintaining cooperation requires individuals to distinguish justified from unjustified defection\cite{sugden1986,kandori1992,leimar2001,panchanathan2003,ohtsuki2006}. However, social norms that sustain cooperation under public assessment may fail when reputations are private and assessments are subject to errors\cite{takahashi2006,uchida2010,hilbe2018,okada2018,perret2021,fujimoto2022,fujimoto2023}. Even when individuals follow the same social norm, they may disagree about whether a particular action is justified. For example, refusing to help a recipient may be considered justified by observers who regard the recipient as bad, but condemned by those who regard the recipient as good\cite{hilbe2018,radzvilavicius2019,okada2020dilemmas}. Such disagreement can extend beyond the individuals involved. Because reputations are both the outcomes of past assessments and the inputs to future assessments, an erroneous judgement of one individual can generate erroneous judgements of others\cite{hilbe2018,uchida2013,fujimoto2022}. Depending on the social norm, this feedback may cause cooperators to lose good standing or allow defectors to acquire favourable reputations\cite{okada2020dilemmas,morsky2024,fujimoto2023}. Reliable reputations thus require both the ability to discriminate between cooperators and defectors and a sufficient degree of agreement among observers\cite{michelmata2024,murase2024synchronization}. Agreement alone is not sufficient, as individuals may agree on reputations that fail to distinguish cooperators from defectors.

The reputational damage caused by errors depends not only on how often individuals misjudge an action, but also on how strongly each judgement changes an established reputation. Previous studies have sought to improve the reliability of reputations through changes in moral assessment\cite{radzvilavicius2019,schmid2021,okada2024}, public institutions\cite{radzvilavicius2021}, gossip\cite{kawakatsu2024} and opinion synchronization\cite{murase2024synchronization,tham2026}. Other studies have shown that quantitative reputations and the aggregation of multiple observations can reduce the effects of erroneous assessments\cite{clark2020,schmid2023,michelmata2024}. These developments point to a distinction between two processes underlying reputation dynamics: assessing an action and revising the reputation of the individual who performed it. The former is the classical question of how goodness should be defined\cite{ohtsuki2004}. The latter concerns how a judgement of an action should change an observer's existing judgement of its actor. In models with instantaneous binary updating, these two processes coincide: a negative assessment immediately assigns a bad reputation, just as a positive assessment immediately restores good standing\cite{ohtsuki2004,ohtsuki2006,hilbe2018}. Yet judging an action as good or bad is not necessarily equivalent to judging the individual as good or bad. A social norm that appropriately evaluates individual actions need not generate reliable reputations if the way these assessments are incorporated into existing judgements is itself inadequate. Understanding the evolutionary consequences of a social norm thus requires considering not only the assessments it produces, but also how these assessments alter reputations over time.

An important feature of reputation updating is that the effect of a new assessment may depend on the reputation already established\cite{mendesiedlecki2013neural,pandula2024}. An isolated negative assessment need not overturn a good reputation built through previous behaviour, just as an isolated positive assessment need not fully restore the reputation of an individual in bad standing. Consistent with this view, recent experimental work on indirect reciprocity suggests that individuals revise reputations gradually rather than replacing them after each observation\cite{yamamoto2025}. Moreover, studies of impression updating have shown that positive and negative behavioural information can have different effects on existing social judgements\cite{reeder1979,skowronski1989,mendesiedlecki2013diagnostic,siegel2018}. Models of indirect reciprocity have likewise allowed upward and downward changes of public reputations to differ\cite{tanabe2013}, and good and bad acts to be weighted unequally\cite{rankin2009}. These observations suggest that reputation loss and recovery need not follow the same updating process. Such asymmetry may have important consequences for indirect reciprocity because different social norms are susceptible to different forms of reputational distortion. Under a norm that frequently assigns negative assessments to conditional cooperators\cite{okada2020dilemmas}, maintaining established good reputations may help preserve the distinction between cooperators and defectors. Under a norm that readily assigns positive assessments to defectors\cite{schmid2021}, however, it may be more important to restrict the recovery of good standing. The effectiveness of asymmetric reputation updating may therefore depend on the social norm that governs behavioural assessment. This raises a fundamental question for the theory of indirect reciprocity: how should individuals revise their reputational judgements, given the social norms they follow, such that reliable reputations and stable cooperation can emerge from imperfect social information?

Here we develop an evolutionary framework of indirect reciprocity that separates the assessment of behaviour from the updating of private reputations. We introduce asymmetric reputational inertia, whereby the loss and recovery of good standing can require different amounts of evidence, and new assessments can affect how individuals are treated before their established reputations are overturned. Combining analytical approaches with individual-based simulations, we show that the effectiveness of reputation updating depends on the social norm under consideration. Different norms benefit from distinct asymmetries that counteract their characteristic forms of reputational distortion, improving both discrimination and agreement. These improvements support stable cooperative outcomes and enlarge their basins of attraction, so that cooperation can be reached from a wider range of initial population compositions. For a given total amount of inertia, its allocation between reputation loss and recovery can lead to contrasting evolutionary outcomes, highlighting the importance of how reputational memory is organized rather than simply how much is retained. We further show how reputational stability and behavioural sensitivity can be jointly regulated, allowing reputations to withstand occasional errors while retaining the reputational consequences of defection. Together, our findings show how the interplay between social norms and reputation updating shapes the collective dynamics of reputation, and thereby the conditions under which indirect reciprocity can sustain cooperation.
